%% ============================================================
\section{Paillier Encryption}
%% ============================================================

Paillier encryption~\cite{Paillier1999} is an additively homomorphic public-key scheme based on the
\emph{decisional composite residuosity} (DCR) assumption. Its design contrasts with textbook RSA: multiplying RSA ciphertexts multiplies plaintext representatives, whereas multiplying Paillier ciphertexts adds their exponents and therefore adds plaintexts modulo $n$.

%% ------------------------------------------------------------
\subsection{Setup and Key Generation}

The key generation algorithm proceeds as follows.

\begin{enumerate}
  \item Sample large primes $p, q$ satisfying $\gcd(pq,\, (p-1)(q-1)) = 1$, and set $n = pq$.
  \item Compute $\lambda \gets \mathrm{lcm}(p-1,\, q-1)$.
  \item Sample $g \gets \Z_{n^2}^*$ uniformly at random.
  \item Compute $\mu \gets \bigl(L(g^\lambda \bmod n^2)\bigr)^{-1} \bmod n$, where
        $L(x) \coloneqq \lfloor (x-1)/n \rfloor$.
        (This requires $L(g^\lambda \bmod n^2)$ to be invertible mod $n$; resample $g$ if not.)
  \item Return $pk \gets (n, g)$ and $sk \gets (\lambda, \mu)$.
\end{enumerate}

\begin{remark}[Simplified generator]
  The common choice $g = 1 + n$ is always valid and simplifies both key generation and
  decryption. Since $(1+n)^m \equiv 1 + mn \pmod{n^2}$ (Lemma~\ref{lem:binomial} below),
  one obtains $L(g^\lambda \bmod n^2) = \lambda \bmod n$, so $\mu = \lambda^{-1} \bmod n$
  is computable directly without sampling. The security level is identical; see
  Remark~\ref{rem:g-security}.
\end{remark}

\paragraph{Encryption.}
$\mathsf{Enc}(pk, m \in \Z_n)$: sample $r \gets \Z_n^*$ and return
\[
  c \;\equiv_{n^2}\; g^m \cdot r^n.
\]

\paragraph{Decryption.}
$\mathsf{Dec}(sk, c)$: compute
\[
  m \;\gets\; L\!\left(c^\lambda \bmod n^2\right) \cdot \mu \bmod n.
\]

\paragraph{Correctness.}
We use the following key lemma.

\begin{lemma}[Binomial structure of $\langle 1+n \rangle$]
\label{lem:binomial}
  For any $w \in \Z_{n^2}^*$, write its unique factorization in the direct product
  $\Z_{n^2}^* \cong \langle 1+n \rangle \times H$ (where $H = \{r^n : r \in \Z_n^*\}$ is
  the subgroup of $n$-th powers) as $w = (1+n)^{[w]} \cdot h$ for $[w] \in \Z_n$ and $h \in H$.
  Then:
  \begin{enumerate}
    \item $(1+n)^k \equiv 1 + kn \pmod{n^2}$ for all $k \in \Z$.
    \item $h^{\lambda} \equiv 1 \pmod{n^2}$ for all $h \in H$.
    \item $L(w^\lambda \bmod n^2) = \lambda [w] \bmod n$.
  \end{enumerate}
\end{lemma}

\begin{proof}
  \textbf{(i)} By the binomial theorem,
  $(1+n)^k = \sum_{j=0}^{k} \binom{k}{j} n^j \equiv 1 + kn \pmod{n^2}$,
  since all terms with $j \ge 2$ are divisible by $n^2$.

  \textbf{(ii)} Any $h \in H$ has the form $h = r^n$ for $r \in \Z_n^*$. The element
  $r^n \in \Z_{n^2}^*$ has order dividing $\phi(n)$ (by Carmichael's theorem applied to
  $n^2$ together with the condition $\gcd(n, \phi(n)) = 1$), and $\lambda \mid \phi(n)$,
  so $h^\lambda = r^{n\lambda} \equiv 1 \pmod{n^2}$.

  \textbf{(iii)} Write $w = (1+n)^{[w]} \cdot h$. Then
  $w^\lambda = (1+n)^{[w]\lambda} \cdot h^\lambda = (1+n)^{[w]\lambda}$
  by (ii). Applying (i): $(1+n)^{[w]\lambda} \equiv 1 + [w]\lambda n \pmod{n^2}$,
  so $L(w^\lambda) = [w]\lambda \bmod n$.
\end{proof}

Correctness of decryption now follows immediately: $c = g^m r^n$ gives
$[c] = m[g]$ (since $[r^n] = 0$), so
\[
  L(c^\lambda) \cdot \mu
  = \lambda[c] \cdot \mu
  = \lambda m [g] \cdot (L(g^\lambda))^{-1}
  = \lambda m [g] \cdot (\lambda [g])^{-1}
  = m \bmod n.
\]

%% ------------------------------------------------------------
\subsection{The DCR Assumption}

\begin{tcolorbox}[colback=blue!4!white,colframe=blue!45!black]
  \textbf{DCR assumption.} Let $n=pq$ be generated as above. For every probabilistic polynomial-time adversary $\mathcal A$,
  \[
  \left|
  \Pr_{r\leftarrow\Z_n^*}[\mathcal A(n,r^n\bmod n^2)=1]
  -
  \Pr_{z\leftarrow\Z_{n^2}^*}[\mathcal A(n,z)=1]
  \right|
  \]
  is negligible in the security parameter.
\end{tcolorbox}

The $n$-th powers in $\Z_{n^2}^*$ form the subgroup $H = \{r^n \bmod n^2 : r \in \Z_n^*\}$
of index $n$ and order $\phi(n)$; the DCR assumption says the $n$ cosets of $H$ are
computationally hidden.

For $g=1+n$, IND-CPA security follows by a direct hybrid. Given a DCR challenge $z$, a reduction chooses a random bit $b$ and gives the adversary
\[
c^*=(1+n)^{m_b}z\bmod n^2.
\]
If $z=r^n$, this is a correctly distributed encryption of $m_b$. If $z$ is uniform in $\Z_{n^2}^*$, multiplication by the fixed element $(1+n)^{m_b}$ preserves uniformity, so $c^*$ is independent of $b$. Any non-negligible IND-CPA advantage therefore distinguishes the two DCR distributions. This is also the usual equivalence between semantic security and IND-CPA for public-key encryption, specialized to Paillier's ciphertext distribution.

\begin{remark}[Security with $g = 1+n$]
\label{rem:g-security}
  Using $g = 1+n$ does not weaken security. The ciphertext $c = (1+mn) \cdot r^n$ is
  uniformly random within the coset $(1+mn) \cdot H$. Under DCR, these cosets are
  computationally indistinguishable, so an adversary cannot determine $m$ from $c$
  regardless of which valid $g$ is used. The coset structure is determined entirely by
  $r^n$, not by the choice of $g$.
\end{remark}

%% ------------------------------------------------------------
\subsection{Additive Homomorphism}

For ciphertexts $c_1 = \mathsf{Enc}(m_1, r_1)$ and $c_2 = \mathsf{Enc}(m_2, r_2)$:
\[
  c_1 \cdot c_2 \bmod n^2
  = g^{m_1} r_1^n \cdot g^{m_2} r_2^n
  = g^{m_1+m_2} (r_1 r_2)^n
  = \mathsf{Enc}(m_1 + m_2 \bmod n,\; r_1 r_2).
\]
Scalar multiplication follows similarly:
\[
  c_1^k \bmod n^2 = g^{km_1}(r_1^k)^n = \mathsf{Enc}(k \cdot m_1 \bmod n,\; r_1^k).
\]
Both operations require only the public key and leave the randomness well-distributed.

%% ============================================================
\section{Threshold Paillier: Construction Sketch}
%% ============================================================

Threshold Paillier distributes decryption authority among $\ell$ parties so that a qualified set can jointly decrypt, while a coalition below the threshold cannot decrypt on its own. The presentation below isolates the reconstruction algebra used by standard threshold variants; complete schemes also specify modulus generation, verifiable shares, and a security proof \cite{FouquePoupardStern2000,DamgardJurik2001}.

%% ------------------------------------------------------------
\subsection{Motivation}

Threshold decryption removes a single decryption authority. It is useful when encrypted data are aggregated publicly but opening the result must require agreement, as in electronic voting, sealed-bid auctions, custody systems, and multiparty computation. Homomorphic evaluation can proceed under the public key, while policy is enforced only at decryption time.

%% ------------------------------------------------------------
\subsection{Key Generation and Share Distribution}

We describe the \emph{trusted dealer} model; see Remark~\ref{rem:dkg} for the DKG variant.

\begin{enumerate}
  \item Run Paillier key generation to obtain $pk = (n, g)$ and $sk = (\lambda, \mu)$.
    Publish $pk$.
  \item Sample a random degree-$(t-1)$ polynomial $f \in \Z_{\phi(n)n}[X]$ with
    $f(0) = \phi(n)$ (so $sk = \phi(n)$).
  \item For each party $i \in [\ell]$, set $sk_i \gets f(i) \bmod \phi(n)n$ and send $sk_i$
    privately to party $i$.
  \item Erase $f$ and $sk$.
\end{enumerate}

\begin{remark}[Composite-ring sharing]
This interpolation takes place over the composite ring $\Z_{\phi(n)n}$, not a field. It is therefore not ordinary Shamir sharing, and perfect privacy cannot be inferred from the standard field proof. Threshold Paillier constructions choose the sharing modulus and secret congruences together so that the required integer interpolation is well defined, and establish privacy by reduction to the underlying encryption assumption. The formulas here should be read as the reconstruction layer of that construction, not as a stand-alone generic secret-sharing scheme.
\end{remark}

Set $\Delta \coloneqq \ell!$. For any subset $S \subseteq [\ell]$ with $|S| = t$ and any
$i \in S$, define the \emph{normalized Lagrange coefficient}
\[
  \Lambda_i(0) \;\coloneqq\; \Delta \cdot \lambda_i^S(0)
  \;=\; \Delta \prod_{j \in S \setminus \{i\}} \frac{-j}{i-j} \;\in\; \Z,
\]
which is always an integer since all denominators divide $\Delta = \ell!$.

\begin{remark}[DKG variant]
\label{rem:dkg}
In a distributed-key-generation variant, the parties jointly generate the RSA modulus and compatible decryption shares without reconstructing $p$, $q$, or the complete secret key at any one party. This is a separate protocol problem and is substantially more involved than distributing shares from a trusted dealer. The construction sketch in this chapter assumes the dealer model.
\end{remark}

%% ------------------------------------------------------------
\subsection{Partial Decryption and Combining}

Given a ciphertext $c$ and a set $S$ of $t$ participating parties, decryption proceeds as
follows.

\paragraph{Partial decryption.}
Each party $i \in S$ computes its \emph{decryption share}:
\[
  c_i \;\equiv_{n^2}\; c^{2\Delta \cdot sk_i}.
\]
(The factor of $2$ ensures that the relevant values lie in the quadratic-residue subgroup used by the threshold construction \cite{DamgardJurik2001}.)

\paragraph{Combining.}
The combiner computes:
\[
  \hat{m} \;\equiv_{n^2}\; \prod_{i \in S} c_i^{2\Lambda_i(0)}.
\]
Since $\sum_{i \in S} sk_i \cdot \Lambda_i(0) = \Delta \cdot sk$ by Lagrange interpolation,
the product evaluates to $c^{4\Delta^2 \cdot sk}$. Applying $L$ and scaling by the
appropriate precomputed constant recovers $m$.

\paragraph{Why integer exponents suffice.}
The Lagrange coefficients $\lambda_i^S(0)$ are rational in general. Multiplying by
$\Delta = \ell!$ clears all denominators (since differences $i - j$ for $i,j \in [\ell]$
divide $\ell!$), making $\Lambda_i(0) \in \Z$. This avoids the need to compute modular
inverses in $\Z_{n^2}^*$, which would require knowing $\phi(n^2)$ and hence the
factorization of $n$ --- precisely the secret we are trying to protect.

%% ------------------------------------------------------------
\subsection{Semi-Honest Security and Simulation}

The scheme above is secure against \emph{semi-honest} adversaries who corrupt up to $t-1$
parties and follow the protocol honestly but try to learn the plaintext from their views.
Security is proved by a reduction to plain Paillier IND-CPA security.

\paragraph{The reduction.}
The reduction is given a plain Paillier public key $N$ and access to a CPA challenge oracle.
It must simulate the entire threshold decryption game for an adversary $A$ who corrupts a
set $C \subseteq [\ell]$ with $|C| \leq t-1$.

\begin{enumerate}
  \item \textbf{Setup.} Sample $\Theta \leftarrow \Z_N$ and set $pk = (N, \Theta)$.
    For each corrupted party $i \in C$, sample $sk_i \leftarrow \Z_{N^2}$ uniformly and
    hand these to $A$. The reduction does not know the real $sk = \phi(N)$.

  \item \textbf{Decryption queries.} When $A$ asks for a decryption of message $m$, the
    reduction:
    \begin{enumerate}[(a)]
      \item Samples $r \leftarrow \Z_N$ and computes $c \equiv_{N^2} (1+mN) \cdot r^N$
        (possible using only $pk = N$).
      \item For each honest party $j \notin C$, computes the simulated share:
        \[
          c_j \;\equiv_{N^2}\;
          (1 + Nm\Theta)^{\Lambda_0(j)} \cdot \prod_{i \in C} c^{sk_i \cdot \Lambda_i(j)},
        \]
        where $\Lambda_0(j) = \Delta \prod_{k \in C} \frac{j-k}{0-k}$ and
        $\Lambda_i(j) = \Delta \cdot \frac{j}{i} \prod_{k \in C \setminus \{i\}} \frac{j-k}{i-k}$
        are normalized Lagrange coefficients for evaluating at $j$ from the set $C \cup \{0\}$.
    \end{enumerate}
    The key insight is that by Lemma~\ref{lem:binomial},
    $c^{sk} = (1+mN)^{\phi(N)} \equiv 1 + mN\phi(N) \equiv 1 + mN\Theta \pmod{N^2}$
    (since $\Theta = \phi(N) \bmod N$ in the real scheme, and here $\Theta$ plays this role).
    So the simulator expresses the honest party's contribution via $m$ and $\Theta$ alone,
    without needing $sk$.

  \item \textbf{Challenge.} $A$ sends $(m_0, m_1) \in \Z_N^2$. The reduction forwards the pair
    \[
      (\Theta\Delta^2 m_0 \bmod N,\ \Theta\Delta^2 m_1 \bmod N)
    \]
    to the plain Paillier CPA oracle and receives $c_b^*$. It forwards $c_b^*$ to $A$ as the challenge ciphertext,
    along with simulated shares computed as above. The messages are scaled by $\Theta\Delta^2$
    because the simulation introduces these extra factors: the combiner ultimately recovers
    $m\Delta^2\Theta \bmod N$ rather than $m$, so the challenge must be embedded accordingly.

  \item \textbf{Output.} $A$ outputs a guess $b'$; the reduction forwards it to the CPA
    challenger. Any advantage of $A$ translates directly to an advantage against plain
    Paillier IND-CPA.
\end{enumerate}

\begin{remark}[Why the simulation works]
  The simulated shares $c_j$ are indistinguishable from real ones because the honest
  party's share is determined by Lagrange interpolation --- from $A$'s view (which includes
  only the corrupted shares and the output), there is exactly one consistent value for each
  honest share, and the simulator produces exactly that value. No information about $sk$
  leaks beyond what is already implied by the corrupted shares and $m$.
\end{remark}

%% ------------------------------------------------------------
\subsection{Achieving Malicious Security via Verifiable Partial Decryptions}

The semi-honest scheme above fails against \emph{malicious} adversaries: a corrupt party
can submit an arbitrary $c_i' \neq c^{2\Delta \cdot sk_i}$ as its decryption share,
causing the combiner to reconstruct a bogus plaintext with no way to detect the fault.
To achieve malicious security, each partial decryption must come with a proof of
correctness.

\paragraph{Intuition.}
At setup the dealer publishes, alongside each share $sk_i$, a \emph{verification key}
$pk_i = v^{\Delta \cdot sk_i}$, where $v$ is a public generator of the subgroup
$\mathbb{Q}_{N^2}$ of quadratic residues modulo $N^2$.  When party $i$ submits
$c_i \equiv_{N^2} c^{2\Delta \cdot sk_i}$, any combiner can check consistency by
verifying that $c_i^2$ and $pk_i$ share the same discrete logarithm --- namely
$\Delta \cdot sk_i$ --- with respect to the bases $c^{4\Delta}$ and $v^{\Delta}$
respectively.  This is a \emph{proof of discrete-log equality (DLEQ)}: a
zero-knowledge argument that
\[
  \log_{c^{4\Delta}} c_i^2 \;=\; \log_{v^{\Delta}} pk_i \;=\; sk_i.
\]
A forged share $c_i' \neq c^{2\Delta \cdot sk_i}$ would require producing a valid DLEQ
proof for an incorrect relation, which is computationally infeasible under the
discrete-logarithm assumption in $\mathbb{Q}_{N^2}$.

\paragraph{Why safe primes and a special $g$?}
Standard DLEQ / Chaum--Pedersen proofs require working in a group of \emph{known} prime
order.  In $\mathbb{Z}_{N^2}^*$ the group order $\phi(N^2)$ is unknown (knowing it
factors $N$).  Two modifications address this:
\begin{enumerate}
  \item \textbf{Safe primes.} Choose $p = 2p'+1$ and $q = 2q'+1$ with $p', q'$ prime.
    Then $\mathbb{Q}_{N^2}$ has order $N p'q'=N\phi(N)/4$. Even though this order is unknown to the prover and verifier, the group has
    no small subgroups, which is the key property needed for the soundness argument.
  \item \textbf{Special generator.} Pick $g \equiv_{N^2} (1+N)^a b^N$ for $b \in \Z_N^*$
    and set $v = g^{2\alpha}$ for random $\alpha \in \Z_N^*$.  This ensures $v$ generates
    $\mathbb{Q}_{N^2}$ and, crucially, the reduction knows $\log_g v = 2\alpha$, which is
    needed to embed the DLEQ proof into a security argument against a plain-Paillier
    CPA adversary.
\end{enumerate}

\paragraph{Extra factors of $2$ and $4$.}
The exponent $2\Delta$ (rather than $\Delta$) in $c_i = c^{2\Delta \cdot sk_i}$ ensures
$c_i \in \mathbb{Q}_{N^2}$ (squaring maps any element into the QR subgroup), so that
DLEQ can be stated entirely within $\mathbb{Q}_{N^2}$.  The DLEQ proof itself introduces
a further squaring, so the combiner ultimately recovers $c^{4\Delta^2 \cdot sk}$ and must
divide out the extra factor of $4\Delta^2$ to obtain $m$.

The share proof follows the Chaum--Pedersen equality-of-discrete-log pattern \cite{ChaumPedersen1992}: the prover commits using fresh exponent randomness, receives a challenge, and returns a linear response that is checked against both base--target pairs. In an unknown-order group the challenge and response ranges must be bounded explicitly; otherwise the ordinary prime-order special-soundness argument does not apply. The exact range proof and soundness analysis are part of the cited threshold-Paillier constructions rather than the elementary derivation above \cite{FouquePoupardStern2000,DamgardJurik2001}.

%% ------------------------------------------------------------
\subsection{Security}

In the threshold IND-CPA experiment, an adversary receives the public key and the internal shares of fewer than $t$ corrupted parties, chooses two messages, and receives an encryption of one of them. Subject to the decryption-oracle restrictions of the chosen model, its advantage in guessing the challenge bit must be negligible. Robust variants additionally require that malformed decryption shares be detected and that any qualified set of valid shares reconstruct the same plaintext. These properties depend on the complete sharing, verification, and setup protocols; they do not follow from the interpolation identity alone \cite{FouquePoupardStern2000,DamgardJurik2001}.
